% QUESTAO
Observe o seguinte polígono
\vspace{-1mm}
\begin{center}
\tikz[scale=2,thick]{
\pgfmathsetmacro{\nn}{9};
\pgfmathsetmacro{\ang}{360/\nn};
\pgfmathsetmacro{\rnd}{100};
\pgfmathsetmacro{\angulo}{270-0.5*\ang+\rnd*rand/\nn};
\coordinate (A) at (\angulo:1);
\pgfmathsetmacro{\angulo}{270+0.5*\ang+\rnd*rand/\nn};
\coordinate (B) at (\angulo:1);
\pgfmathsetmacro{\angulo}{270+1.5*\ang+\rnd*rand/\nn};
\coordinate (C) at (\angulo:1);
\pgfmathsetmacro{\angulo}{270+2.5*\ang+\rnd*rand/\nn};
\coordinate (D) at (\angulo:1);
\pgfmathsetmacro{\angulo}{270+3.5*\ang+\rnd*rand/\nn};
\coordinate (E) at (\angulo:1);
\pgfmathsetmacro{\angulo}{270+4.5*\ang+\rnd*rand/\nn};
\coordinate (F) at (\angulo:1);
\pgfmathsetmacro{\angulo}{270+5.5*\ang+\rnd*rand/\nn};
\coordinate (G) at (\angulo:1);
\pgfmathsetmacro{\angulo}{270+6.5*\ang+\rnd*rand/\nn};
\coordinate (H) at (\angulo:1);
\pgfmathsetmacro{\angulo}{270+7.5*\ang+\rnd*rand/\nn};
\coordinate (I) at (\angulo:1);
\draw (A)--(B)--(C)--(D)--(E)--(F)--(G)--(H)--(I)--cycle;
\fill (E) circle (0.7pt);
}
\end{center}
\vspace{-2mm}
Determine o número de diagonais que partem de um único vértice.

% ALTERNATIVAS 3
$6$ diagonais
$5$ diagonais
$4$ diagonais
$3$ diagonais
$7$ diagonais


